\section[Comparison with the Bloch--Kato scalar package]{Comparison with the Bloch--Kato scalar package}
\label{sec:comparison}

The preceding sections built the five columns separately.  This section
contracts them back to the usual Bloch--Kato scalar package and asks two
questions: is the typed source adequate for that contraction, and what
does the contraction forget?

\begin{remark}[Relation to the standard presentations]\label{rmk:apparatus:five-column-literature-gap}
The fundamental line and its Tamagawa-number formulation go back to
Bloch--Kato~\citep[\S5]{BlochKato1990} and
Fontaine--Perrin-Riou~\citep{FontainePerrinRiou1994}; the formulation of
the leading value in terms of periods and height determinants, modulo
$\Q^\times$, is surveyed by Nekov\'a\v{r}~\citep{Nekovar1994}.  The
five-slot organization used here is a bookkeeping choice made for the
adequacy tests.  We have not compared it systematically with the
literature and make no claim about whether equivalent decompositions
appear there.
\end{remark}

\begin{definition}[Five-column tensor]\label{def:apparatus:five-column-tensor}
The typed five-column tensor is
\[
  T_{\mathrm{5col}}(E)
  =
  \rho_{\mathrm{ht/reg}}^{\R}(E)
  \otimes
  \rho_{\mathrm{an/period}}^{\R}(E)
  \otimes
  \rho_{\mathrm{finite}}^{\R}(E)
  \otimes
  \bigotimes_{p\in S(E)}\rho_p^{\R}(E)
  \otimes
  \rho_{\det}^{\R}(o_{\det}),
\]
whose fifth factor is the orientation scalar of the independent
orientation datum (\Cref{def:apparatus:det-assembly-map}), not the final
assembly $\mathcal A^{\R}(E)$.  Its source family consists of the five
separated logarithmic tangent coordinates
$(x_{\mathrm{ht}},x_{\mathrm{an}},x_{\mathrm{fin}},x_p,x_{\det})$, where
$x_{\det}=\log|\rho_{\det}^{\R}(o_{\det})|$.  The
\emph{orientation-normalized scalar readout} is
\[
  \log\bigl|T_{\mathrm{5col}}(E)\big/\rho_{\det}^{\R}(o_{\det})\bigr|
  =x_{\mathrm{ht}}+x_{\mathrm{an}}+x_{\mathrm{fin}}+x_p ,
\]
the logarithm of the absolute value of the product of the four arithmetic
factors.  It is distinct from the full contraction
$\log|T_{\mathrm{5col}}(E)|$, whose row is $(1,1,1,1,1)$.
\end{definition}

The distinction in the definition matters.  Since $\mathcal A^{\R}(E)$
already contains the four arithmetic factors, using it as the fifth factor
would count each of them twice.  For instance, with arithmetic factors
$2,3,5,7$ and orientation scalar $1$, the final assembly is $210$; the
tensor with the orientation as fifth factor is again $210$, whereas
reusing the final assembly would give $210^2=44100$.  (Lean checks these
values.)

\begin{theorem}[Five-column to Bloch--Kato scalar contraction]\label{thm:apparatus:five-col-bk-comparison}
In the coordinates
$x=(x_{\mathrm{ht}},x_{\mathrm{an}},x_{\mathrm{fin}},x_p,x_{\det})\in\R^5$,
the orientation-normalized scalar readout of
\Cref{def:apparatus:five-column-tensor} is
$\Phi_{\mathrm{5col}\to\mathrm{BK}}(x)=d\cdot x$ with $d=(1,1,1,1,0)$.
With the full typed source $L(x)=x$ and $C=I_5$,
\[
  \Xi_{\mathrm{5col}\to\mathrm{BK}}
  =
  \|d\|^2-d\,I_5\,d^{\mathsf T}
  =
  4-4=0,
\]
so the typed source is adequate for this readout.  The orientation
direction $e_{\det}=(0,0,0,0,1)$ lies in the kernel of
$\Phi_{\mathrm{5col}\to\mathrm{BK}}$: the readout forgets at least that
coordinate.%
\footnote{Lean proves the contraction row, the zero residual and the
kernel vector.  Identifying this coordinate readout with the arithmetic
Bloch--Kato scalar requires the five transports of the column sections
and, in addition, a convention under which that scalar is independent of
the orientation datum.}
\end{theorem}

\begin{proof}
With the identity source, $\KLL=\Kdagger=I_5$, $\KDL=d$ and
$K_{LD}=d^{\mathsf T}$, so the projected energy is
$d\,I_5\,d^{\mathsf T}=dd^{\mathsf T}=4$, which equals $\KDD$; hence the
residual is $0$.  Finally $d\cdot e_{\det}=0$, so changing only
$x_{\det}$ changes the typed coordinate without changing the scalar
readout.
\end{proof}

The two halves of the theorem fit together.  The orientation-normalized
readout is computed correctly from the typed source, so nothing is lost
in that direction; but the typed source carries an orientation coordinate
that this readout does not see, while the full tensor does (changing the
orientation scalar from $1$ to $2$ in the example above changes the
tensor from $210$ to $420$).  To the extent that the arithmetic
Bloch--Kato scalar is orientation-normalized in this sense, which is an
additional convention not supplied by the transports alone, it is a
projection of the typed atlas.
